\documentclass[unicode]{beamer}

% --- [設定] ---
\usepackage{luatexja}
\usepackage[orientation=portrait,size=a0,scale=1.4]{beamerposter}
\usepackage[most]{tcolorbox}
\usepackage{tikz}
\usetheme{default}

\setbeamertemplate{navigation symbols}{}
% ページ全体の基本余白を2cmに設定
\setbeamersize{text margin left=2cm, text margin right=2cm}

% --- [フォントの設定：CloudLaTeX用] ---
\usepackage[haranoaji,deluxe]{luatexja-preset} % 原ノ味フォントを使用するプリセット
\renewcommand{\kanjifamilydefault}{\gtdefault}  % 日本語の標準をゴシック体にする
\renewcommand{\familydefault}{\sfdefault}      % 英数字の標準をサンセリフ（Arial風）にする

% --- [色の設定] ---
% ブロックの見出し（block title）の色をカスタマイズ
\setbeamercolor{block title}{fg=white, bg=cyan!70!black} % 文字を白、背景を濃い水色に
\setbeamercolor{block body}{bg=white}                   % 本文の背景を白に


\begin{document}
\begin{frame}[t, fragile]

  % --- 1. タイトル部分 ---
  \vspace{1cm} % 紙の一番上との間の隙間
  \nointerlineskip
  \begin{tcolorbox}[
      enhanced, sharp corners, boxrule=0pt,
      colback=cyan!20, 
      width=\textwidth,             % 本文の幅（左右余白を除いた幅）に合わせる
      % grow to left/right を消すことで、setbeamersize で設定した余白が活きます
      halign=center, height=10cm
    ]
    {\Huge \bfseries A0日本語ポスター タイトル} \par
    \vspace{1cm}
    {\raggedleft \Large 所属 \hspace{2cm} \par}
    {\raggedleft \Large 著者 \hspace{2cm} \par}
  \end{tcolorbox}

  % --- 2. 中央の罫線（開始位置を調整） ---
  \begin{tikzpicture}[remember picture, overlay]
    \draw [gray!50, line width=3pt] 
      % yshift を -13cm くらいにすると、高さ10cmの箱の下に少し隙間が空きます
      ([yshift=-13cm]current page.north) -- ([yshift=2cm]current page.south);
  \end{tikzpicture}

  \vspace{3cm} % タイトルボックスと本文の間の距離

  % --- 3. 2段組み ---
  \begin{columns}[t] 
    \begin{column}{.48\linewidth}
      \begin{block}{\large ひだり}
        左側のコンテンツ。
      \end{block}
    \end{column}

    \begin{column}{.02\linewidth}
      % 中央の線のスペース
    \end{column}

    \begin{column}{.48\linewidth}
      \begin{block}{\large みぎ}
        右側のコンテンツ。
      \end{block}
    \end{column}
  \end{columns}

\end{frame}
\end{document}
